\subsection{在主理想整环中分解为不可约因子}

我们现在将 VI 第 18 页中关于分解为不可约元素的结果应用于主理想整环。根据命题 2，A 的一个元素 \(p \neq 0\) 是不可约的，当且仅当环 \(\mathrm{A}/(p)\) 是一个域（I，第 115 页，推论 1），即当且仅当同余式 \(ax \equiv b \pmod{p}\) 对每个 \(b \in \mathrm{A}\) 以及每个 \(a \in \mathrm{A}\) （不是 \(p\) 的倍元）在 A 中有解.

定义2. --- 设 A 是一个整环。A 的不可约元素的一个代表系是 A 的不可约元素的一个族 \((p_{\alpha})\) ，使得 A 的每个不可约元素恰好是其中一个 \(p_{\alpha}\) 的相伴元.

定理2. --- 设 A 是一个主理想整环，且设 \((p_{\alpha})\) 为 A 的不可约元素的一个代表系。则A的分式域中每个非零元素x都可以唯一地表示为如下形式

\[
x = u \prod_ {\alpha} p _ {\alpha} ^ {n _ {\alpha}},\tag{1}
\]

其中 u 是 A 的一个可逆元，且诸 \(n_{\alpha}\) 都是整数，除有限多个外均为零。要使 x 属于 A，必要且充分的条件是所有 \(n_{\alpha}\) 都是正的。

我们将使用关于分解为不可约元素之和的定理（VI，第 18 页，定理 2），上述陈述只是该定理的一种翻译。由于 \(\mathcal{P}^*\) 是一个格序群，为了验证该定理的假设确实成立，我们只需证明A的主理想的每个非空集合都包含一个极大元；而这由下面的引理得出：

引理1. --- 设 A 是一个环，使得 A 的每个左理想都是有限生成的。那么 A 的左理想组成的、按包含关系排序的每个非空集合 \(\Phi\) 都具有一个极大元。

由佐恩引理（集合论，III，第 154 页，定理 2）可知，只需证明 \(\Phi\) 是归纳的。现在如果 \((\mathfrak{a}_{\lambda})\) 是 \(\Phi\) 中元素的一个全序族，那么并集 \(\mathfrak{a}\) （诸理想 \(\mathfrak{a}_{\lambda}\) 的）是 \(\mathbf{A}\) 的左理想，因此具有有限生成系 \((a_i)_{1 \leqslant i \leqslant n}\) 。由于每个 \(a_i\) 属于某个理想 \(\mathfrak{a}_{\lambda_i}\) ，并且由于族 \((\mathfrak{a}_{\lambda})\) 是全序的，诸 \(a_i\) 都属于诸理想 \(\mathfrak{a}_{\lambda_i}\) 中最大的一个，设为 \(\mathfrak{a}_{\mu}\) 。于是 \(\mathfrak{a} = \mathfrak{a}_{\mu}\) 属于 \(\Phi\) ，后者因此确实是一个归纳集。

稍后我们将研究那些环 B，称为诺特环，使得 B 的每个非空理想集都包含一个极大元。

注记。 --- 族 \((u, (n_{\alpha}))\) 称为 x 分解为不可约因子的分解；滥用语言地，我们也称公式 (1) 是 x 分解为不可约因子的分解。如果 \(x = u \prod_{\alpha} p_{\alpha}^{n_{\alpha}}\) 与 \(y = v \prod_{\alpha} p_{\alpha}^{m_{\alpha}}\) 是

是 x 和 y 分解为不可约因子的分解，那么 x 整除 y 的一个必要且充分条件是 \(n_{\alpha} \leqslant m_{\alpha}\) 对所有 \(\alpha\) ；由此我们推导出以下公式

\[
\operatorname * {g c d} (x, y) = \prod_ {\alpha} p _ {\alpha} ^ {\inf (n _ {\alpha}, m _ {\alpha})}\tag{2}
\]

\[
\operatorname{lcm} (x, y) = \prod_ {\alpha} p _ {\alpha} ^ {\sup \left(n _ {\alpha}, m _ {\alpha}\right)}.\tag{3}
\]

定理2所表达的性质对于比主理想整环更广泛的一类环也成立；我们稍后将把它们作为唯一分解整环来研究；并且我们将看到，在任意多个未定元上的多项式环和形式幂级数环都是唯一分解整环（交换代数，第七章，第3节）。

